\chapter*{Conclusion}
\addcontentsline{toc}{chapter}{Conclusion}
\label{ch:conclusion}
\index{revision checklist}

You now have a map of the terrain.  Part~\ref{part:startup} gave you
the startup lifecycle, GC tuning, and a diagnostic playbook.
Part~\ref{part:packages} explained the package-management and loading
model, the layer where most configuration bugs hide.
Part~\ref{part:lang} traced the path from file extension to
fully-powered editing via tree-sitter, Eglot, and DAP.
Part~\ref{part:ui} covered the keymap stack, the two completion axes,
and the buffer/window/frame display model.
Part~\ref{part:org} addressed org-mode's loading cost, async export,
babel registration, and the credential infrastructure underneath.
Part~\ref{part:idioms} distilled everything into the five recurring
patterns that account for most real-world configuration code.

No single reading can internalise all of it.  The \textbf{Mental
model} summaries at the end of each chapter are designed as revision
checklists; skim them periodically when you return to a topic.  The
appendices (variable table, hook timing, keybinding index, and the
deep-dive supplements on GC, LSP, and straight.el) serve as
quick-reference material you can consult without re-reading full
chapters.

Emacs rewards curiosity.  When something surprises you (a
mysterious pause, an unexpected keybinding, a package that refuses to
load), reach for the diagnostic tools described in
Chapter~\ref{ch:profiling} and the playbook in
Chapter~\ref{ch:diagnostic-playbook}.  The answer is always in the
source; this book has shown you where to look.
